\clearpage\begingroup\SourceLeaf{432b}\ArithmeticHead{101}{R}{DIE ALLGEMEINE LEHRE VON DEN KRUMMEN FLÄCHEN.}
\renewcommand{\thefootnote}{\arabic{footnote})}\setcounter{footnote}{0}
\begin{center}
\textbf{31.}\par\smallskip
\textls[100]{Die Entstehung der Disquisitiones generales circa superficies}\par
\textls[100]{curvas} (1826—1827).
\end{center}
\noindent\hspace*{\parindent}Nur nach langem, hartem Ringen hat \textsc{Gauss} das Ziel erreicht, das er
sich Ende 1825 gestellt hatte, die Lehre von den krummen Flächen in voller
Allgemeinheit zu begründen. Am 19. Februar 1826 schreibt er an \textsc{Olbers}:
»Ich wüsste kaum eine Periode meines Lebens, wo ich bei so angestrengter
Arbeit wie in diesem Winter doch verhältnismässig so wenig reinen Gewinn
geerntet hätte. Ich habe \textls[100]{viel}, viel Schönes herausgebracht, aber dagegen
sind meine Bemühungen über anderes oft \textls[100]{Monate} lang fruchtlos gewesen
(Br. G.-O. 2, S. 438). Und am 2. April 1826: »Meine theoretischen Arbeiten
lassen bei ihrem so \textls[100]{sehr} grossen Umfange leider noch viele Lücken; am
leichtesten wäre mir geholfen, wenn ich mir erlaubte, mit der Bekanntmachung
meiner Messungen zwar alle meine Rechnungseinrichtungen zu verbinden, aber
deren \textls[100]{Ableitungen} aus ihren höhern Gründen für ein ganz getrenntes Werk
für glücklichere zukünftige Zeiten aufsparte. Dann wäre nirgends ein Anstoss.
Vors erste werde ich die scharfe Ausgleichung meiner 32 Punkte, die 51 Dreiecke
und 146 Richtungen liefern, vornehmen« (W. IX, S. 376). Es sei hierzu
bemerkt, dass die Arbeiten im Felde im August 1825 beendet waren, und es
sich lediglich um den Abschluss der Rechnungen handelte, sodass \textsc{Gauss} für
seine theoretischen Arbeiten Zeit gewann.
\clearpage\SourceLeaf{433}\ArithmeticHead{102}{L}{STÄCKEL, GAUSS ALS GEOMETER.}\setcounter{footnote}{0}
\indent Im Herbst 1826 scheint \textsc{Gauss} durchgedrungen zu sein. Er berichtet
am 20. November an \textsc{Bessel}: »Die Verarbeitung der Materialien zu dem beabsichtigten
Werke über meine Messungen kostet mich viele Zeit. Meine
Hauptdreiecke, 33 Punkte befassend, sind zwar längst fertig berechnet, aber
die Berechnung der vielen geschnittenen Nebenpunkte .\ .\ . macht viel Arbeit.
.\ .\ . Noch viel mehr Verlegenheit macht mir der weit ausgedehntere
theoretische Teil, der so vielfach in andere Teile der Mathematik eingreift.
Ich sehe hier kein anderes Mittel, als mehrere grosse Hauptpartien von dem
Werke abzutrennen, damit sie selbständig und in gehöriger Ausführlichkeit
entwickelt werden können. Gewissermassen habe ich damit schon in meiner
Schrift über die Abbildung der Flächen unter Erhaltung der Ähnlichkeit der
kleinsten Teile den Anfang gemacht; eine zweite Abhandlung, die ich vor
ein paar Monaten der Königlichen Sozietät übergeben habe und die hoffentlich
bald gedruckt werden wird, enthält die Grundsätze und Methoden zur
Ausgleichung der Messungen\footnote{\textsc{C. F. Gauss}, \textit{Supplementum theoriae combinationis observationum erroribus minimis obnoxiae}, vorgelegt den 16. September 1826, W. IV, S. 55.}. .\ .\ . Vielleicht werde ich zunächst erst noch
eine dritte Abhandlung ausarbeiten, die mancherlei neue Lehrsätze über
krumme Flächen, kürzeste Linien, Darstellung krummer Flächen in der Ebene
usw. entwickeln wird. Hätten alle diese Gegenstände in mein projektiertes
Werk aufgenommen werden sollen, so hätte ich entweder manches ungründlich
abfertigen oder dem Werk ein sehr buntscheckiges Ansehen geben müssen«
(W. IX, S. 362).
\par
Mit der Ausarbeitung der dritten Abhandlung hat \textsc{Gauss} bald darauf begonnen.
Nach dem Briefe an \textsc{Olbers} vom 14. Januar 1827 (Br. G.-O. 2,
S. 467) war er damals »schon ziemlich damit vorgerückt«, und am 1. März 1827
schreibt er jenem, die Abhandlung sei vollendet, er werde sie jedoch der
Sozietät noch nicht übergeben, da doch auf die Ostermesse kein Band der
Denkschriften herauskomme (W. IX, S. 377). In der Tat sind die \textit{Disquisitiones
generales circa superficies curvas} der Göttinger Gesellschaft der Wissenschaften
erst am 8. Oktober 1827 vorgelegt und in den Band VI der Commentationes
recentiores vom Jahre 1828 aufgenommen worden (W. IV, S. 217).
Vorher war in den Göttinger Gelehrten Anzeigen vom 5. November 1827 eine
ausführliche Selbstanzeige erschienen (W. IV, S. 341—347).
\clearpage\SourceLeaf{434}\ArithmeticHead{103}{R}{DIE ALLGEMEINE LEHRE VON DEN KRUMMEN FLÄCHEN.}\setcounter{footnote}{0}
\indent Von den Neuen allgemeinen Untersuchungen unterscheiden sich die \textit{Disquisitiones
generales} hauptsächlich in zwei Punkten: sie enthalten erstens den
Ausdruck für das Krümmungsmass bei beliebiger Wahl der bestimmenden
Veränderlichen \(p,q\) und zweitens die Verallgemeinerung des \textsc{Legendreschen}
Theorems von der Kugel auf beliebige Flächen.
\par
Die allgemeine Formel für das Krümmungsmass hat sich \textsc{Gauss}, wie
schon angedeutet wurde, im Laufe des Jahres 1826 erarbeitet. Die im Nachlass
befindlichen Aufzeichnungen gestatten es hier, einmal einen vollständigen
Einblick in die Entstehung seiner Gedanken zu gewinnen. Da \textsc{Gauss} dabei
an schon vorliegende Untersuchungen über die Abwicklung krummer Flächen
anknüpft, wird es angebracht sein, einen kurzen geschichtlichen Überblick
vorauszuschicken\footnote{Ausführliche Angaben findet man bei \textsc{P. Stäckel}, \textit{Bemerkungen zur Geschichte der geodätischen Linien}, Leipziger Berichte, 1893, S. 452—455.}.
\par
Die Abwicklung von Zylindern und Kegeln auf die Ebene war im 18.
Jahrhundert wiederholt betrachtet und zur Lösung von Aufgaben benutzt
worden. \textsc{Euler} hatte dann (1770) nach den krummen Flächen gefragt, die
sich überhaupt auf eine Ebene abwickeln lassen, und war, indem er der Anschauung
entnahm, dass die gesuchten Flächen gradlinig sein müssen, zu ihrer
allgemeinen Darstellung gelangt\footnote{\textsc{L. Euler}, \textit{De solidis, quorum superficiem in planum explicare licet}, Novi Comment. Petrop. 16 (1771), 1772, S. 3; vorgelegt am 5. März 1770.}. Wie eine erst im Jahre 1862, also nach dem
Tode von \textsc{Gauss} aus \textsc{Eulers} Nachlass abgedruckte Notiz\footnote{\textsc{L. Euler}, \textit{Opera postuma}, St. Petersburg 1862, t. I, S. 494—496.} zeigt, ist dieser um
dieselbe Zeit zu dem Problem gelangt, »invenire duas superficies, quarum
alteram in alteram transformare licet, ita ut in utraque singula puncta homologa
easdem inter se teneant distantias«, und er hat dafür genau die Gleichungen
angesetzt, die man im art. 12 der \textit{Disq. gen.} findet. Es ist ihm auch
gelungen, ihre Integration für die Biegung von Kegeln in Kegel durchzuführen,
und er hat zum Schluss die Frage nach den Biegungen von Stücken einer
Kugelfläche aufgeworfen.
\par
Unabhängig von \textsc{Euler} hatte \textsc{Monge} Untersuchungen über die auf die
Ebene abwickelbaren Flächen angestellt. Er hat sie, durch \textsc{Eulers} Abhandlung
vom Jahre 1771 veranlasst, in einer zweiten Arbeit weiter geführt; in
\clearpage\SourceLeaf{435}\ArithmeticHead{104}{L}{STÄCKEL, GAUSS ALS GEOMETER.}\setcounter{footnote}{0}
\noindent dieser findet sich auch die bekannte partielle Differentialgleichung zweiter
Ordnung für die auf die Ebene abwickelbaren Flächen\footnote{\textsc{G. Monge}, \textit{Sur les développées, les rayons de courbure et les différents genres d’inflexions des courbes à double courbure}, Mém. sav. étr. t. 10, Paris 1785, S. 511 (eingereicht 1771); \textit{Sur les propriétés de plusieurs genres de surfaces courbes, particulièrement sur celles des surfaces développables, avec une application à la théorie des ombres et pénombres}, Mém. sav. étr. t. 9, Paris 1780, S. 382 (eingereicht 1775); vgl. auch \textsc{J. Meusnier}, \textit{Sur la courbure des surfaces}, Mém. sav. étr. t. 10, Paris 1785, S. 509 (vorgelegt 1776).}.
\par
Wie \textsc{Gauss} zu dem allgemeinen Begriff der Biegung krummer Flächen
gelangt ist, wissen wir nicht. Es ist nicht wahrscheinlich, dass die bereits
erwähnten Arbeiten über die Gestalt elastischer Flächen, an denen sich ausser
\textsc{Poisson} (1814) auch \textsc{Lagrange} (1811) und \textsc{Sophie Germain} (1815) beteiligt
hatten\footnote{\textsc{J. L. Lagrange}, \textit{Mécanique analytique}, 2. éd., t. I, Paris 1812, Statique, sect. V, Chap. 3, § II; \textit{De l’équilibre d’un fil ou d’une surface flexible et au même temps extensible et contractible}, Oeuvres, t. 11, S. 156; \textsc{S. Germain}, \textit{Recherches sur la théorie des surfaces élastiques}, Paris 1820 (verfasst 1815). Diese Untersuchungen waren veranlasst durch \textsc{Chladnis} Entdeckungen über die Klangfiguren (\textit{Akustik} Leipzig 1802).}, auf ihn Einfluss gehabt haben. Dagegen sind ihm die schon früher
veröffentlichten Abhandlungen von \textsc{Euler} und \textsc{Monge} bekannt gewesen. Die
auf die Ebene »abwicklungsfähigen Flächen« hat er bereits in der Aufzeichnung
vom Dezember 1822 (W. VIII, S. 382—384) nach der Seite des Krümmungsmasses
betrachtet, und in den \textit{Neuen allgemeinen Untersuchungen} (art. 16) bemerkt
er, aus dem Satze von der Erhaltung des Krümmungsmasses folge der
wichtige, aber bis jetzt nicht mit der wünschenswerten Evidenz abgeleitete
Lehrsatz, dass bei jenen Flächen das Krümmungsmass verschwindet, und damit
sei erst bewiesen, dass sie der bekannten Differentialgleichung genügen (vgl.
auch W. VIII, S. 437 und 444).
\par
Wie wir sahen, hatte \textsc{Gauss} bei zwei besonderen Formen des Linienelementes
das Krümmungsmass durch den darin auftretenden Koeffizienten
und dessen erste und zweite partielle Ableitungen ausdrücken können. Er
wusste, dass das Krümmungsmass bei den Biegungen erhalten bleibt, folglich
musste bei der allgemeinen Form des Linienelementes das Krümmungsmass
ebenfalls durch die darin auftretenden Koeffizienten und deren partielle Ableitungen
darstellbar sein. Allein die Rechnungen, die dort zum Ziel geführt
hatten, liessen sich nicht ohne Weiteres auf den Fall beliebiger bestimmender
Grössen übertragen; hierauf beziehen sich wohl die Klagen über die Unfruchtbarkeit
\clearpage\SourceLeaf{436}\ArithmeticHead{105}{R}{DIE ALLGEMEINE LEHRE VON DEN KRUMMEN FLÄCHEN.}\setcounter{footnote}{0}
\noindent langer Bemühungen in dem schon angeführten Briefe an \textsc{Olbers}
vom 19. Februar 1826.
\par
Im Sommer oder Herbst des Jahres kam \textsc{Gauss} auf den Gedanken, die
auf die Ebene abwickelbaren Flächen heranzuziehen. Diese sind einerseits
dadurch gekennzeichnet, dass das Krümmungsmass verschwindet, andererseits
aber dadurch, dass für sie
\[
E\,dp^2+2F\,dp\,dq+G\,dq^2=dt^2+du^2,
\]
\noindent das heisst gleich dem Produkt der beiden vollständigen Differentiale \(d\lambda=dt+idu\)
und \(d\mu=dt+idu\) ist. \textsc{Gauss} verschaffte sich jetzt (W. VIII, S. 446,
Handbuch 16 Bb, S. 114) die Bedingungsgleichung dafür, dass
\[
E\,dp^2+2F\,dp\,dq+G\,dq^2=d\lambda\,d\mu
\]
\noindent wird. Es ergab sich als linke Seite ein Ausdruck, der aus den Koeffizienten
\(E,F,G\) und deren ersten und zweiten partiellen Ableitungen nach \(p\) und \(q\)
zusammengesetzt ist, und man durfte vermuten, dass er sich vom Krümmungsmass
nur um einen unwesentlichen Faktor unterscheidet.
\par
Damit war \textsc{Gauss} in den Besitz des Zählers gelangt, der bei dem allgemeinen
Ausdruck für das Krümmungsmass auftritt, und nachdem er so das
Ergebnis kannte, glückte es ihm auch, die unmittelbare Ableitung der Formel
zu finden, die im art. 11 der \textit{Disq. gen.} angegeben wird. In der Tat stehen
die Rechnungen über das »Krümmungsmass der Flächen bei allgemeinem Ausdruck
derselben« im Handbuch 16 Bb, S. 128—131, also einige Seiten hinter
der vorher erwähnten Aufzeichnung über die auf die Ebene abwickelbaren
Flächen.
\par
Von dem höheren Standpunkte aus betrachtet, den \textsc{Gauss} jetzt gewonnen
hatte, verlor der ursprüngliche Beweis für die Erhaltung des Krümmungsmasses
in seinen Augen an Wert, ja noch mehr, der Satz von der Winkelsumme
des geodätischen Dreiecks, der dafür den Ausgangspunkt gebildet hatte,
bekam jetzt seine Stelle als eine Folgerung aus dem Hauptsatze von dem
Krümmungsmass, wenn man ihn nämlich auf geodätische Polarkoordinaten
anwandte (\textit{Disq. gen.} art. 20).
\par
Die Verallgemeinerung des \textsc{Legendreschen} Theorems, zu der wir uns
nunmehr wenden, hätte \textsc{Gauss} schon in die \textit{Neuen allgemeinen Untersuchungen}
\par\smallskip\noindent{\small X\textsubscript{2} Abh. 4.\hfill 14}\par
\clearpage\SourceLeaf{437a}\ArithmeticHead{106}{L}{STÄCKEL, GAUSS ALS GEOMETER.}\setcounter{footnote}{0}
\noindent aufnehmen können, denn es war ihm im November 1825 gelungen, sie auf
die elegante Gestalt zu bringen, die er in den artt. 25 bis 28 der \textit{Disq. gen.}
mitteilt, und er würde es sicherlich getan haben, wenn er nicht Ende 1826
die Arbeit an dem Entwurf abgebrochen hätte. Bei jener Verallgemeinerung
wird die Lehre von den kürzesten Linien mit der Lehre vom Krümmungsmass
verbunden, auf die sich die beiden Hauptabschnitte jener Abhandlung beziehen,
und so erscheint der Satz von der Zurückführung kleiner geodätischer Dreiecke
auf ebene Dreiecke als die Krönung des Gebäudes der allgemeinen Lehre
von den krummen Flächen. Zugleich aber bildet er in echt \textsc{Gaussscher} Art
den Übergang zu den Anwendungen. \textsc{Gauss} hat sich hierüber in dem Briefe
an \textsc{Olbers} vom 1. März 1827 folgendermassen ausgesprochen: »Jene Abhandlung
enthält zur \textls[100]{unmittelbaren} Benutzung in meinem künftigen Werk über
die Messung eigentlich nur ein paar Sätze, nämlich:
\par
1) was zur Berechnung des Exzesses der Summe der 3 Winkel über \(180^0\)
in einem Dreiecke auf einer \textls[100]{nicht} sphärischen Fläche, wo die Seiten kürzeste
Linien sind, erforderlich ist,
\par
2) wie in diesem Fall der Exzess \textls[100]{ungleich} verteilt werden muss, damit
die Sinus den Seiten gegenüber proportional werden.
\par
In praktischer Rücksicht ist dies zwar ganz unwichtig, weil in der Tat
bei den grössten Dreiecken, die sich auf der Erde messen lassen, die Ungleichheit
in der Verteilung unmerklich wird; aber die Würde der Wissenschaft
erfordert doch, dass man die Natur dieser Ungleichheit klar begreife. Und
so kann man allerdings hier, wie öfters, ausrufen: Tantae molis erat! um
dahin zu gelangen. — Wichtiger aber als die Auflösung dieser 2 Aufgaben
ist es, dass die Abhandlung mehrere allgemeine Prinzipien begründet, aus
denen künftig, in einer speziellern Untersuchung, die Auflösung von einer
Menge wichtiger Aufgaben abgeleitet werden kann« (W. IX, S. 378).
\clearpage\endgroup
